% УВОД, без номер. Числа тук НЕ се пишат: те се появяват едва в Раздел IV.
\section*{УВОД}
\label{sec:introduction}
\addcontentsline{toc}{section}{УВОД}

Практически важните задачи за комбинаторна оптимизация са изчислително трудни, тоест
точното им решаване е приложимо само за малки входове. Вместо това инженерната практика
използва \term{метаевристики}: общи стратегии за търсене, които се отказват от гаранцията
за оптималност в замяна на контрол върху изразходваното време~\cite{blum2003metaheuristics}.
Формата им, много слабо свързани опити вместо една дълга зависима верига, ги прави удобни
за паралелно изпълнение.

Лесното паралелизиране обаче е подвеждащо. Класификацията
в~\cite{crainic2010parallel} показва, че отделните схеми, независим многократен старт,
споделяне на общо най-добро решение и островен модел, се различават не само по
постигнатото ускорение, а и по това дали решават същата задача: кооперативните схеми
променят траекторията на търсенето, тоест променят и качеството на решението. Оттук и
проблемът на проекта. Твърдението „реализацията е паралелна и затова е по-бърза“ не е
проверимо, докато не бъдат фиксирани едновременно наборът еталонни инстанции, критерият
за спиране и начинът за сравняване на качеството. Без тези три условия измереното
ускорение може да идва от по-слаба работа на нишка, а не от по-добро използване на
машината.

\textbf{Целта} е да се проектира и реализира паралелна рамка за метаевристична
оптимизация, в която задачата, метаевристиката и схемата за паралелно изпълнение стоят
зад общ договор и се менят независимо, и да се направи количествена оценка на
производителността и на качеството при контролирани условия. Задачите са четири:
обосновка на проектните решения; реализация на четири метаевристики и три паралелни
схеми с внимание към цената на синхронизацията; методика за измерване, която фиксира
инстанциите, критерия за спиране, броя повторения и обобщаването; и измерване на
ускорението, на ефективността, на качеството и на разпределението на времето.

\textbf{Обхватът} е рамката, метаевристиките в нея и експерименталната оценка. Извън
обхвата остават разпределеното изпълнение върху няколко машини, изпълнението върху
графични ускорители и точните методи, използвани само като източник на известни
оптимални стойности.

Проектът не предлага нова метаевристика и не твърди подобрение спрямо публикувани
резултати. Той предлага условия, при които сравнението между вече известни методи е
проверимо, и показва с измерване докъде тези условия стигат.
